%=====================================================================
\chapter{Research Contributions}
%=====================================================================

This chapter states what has been contributed against each research objective and
identifies the published work in which that contribution appears. Each publication is
described once, at the point where its contribution is actually made; later chapters refer
to it by its label \pubcite{P1}--\pubcite{P6} rather than restating it, and the
methodology, implementation and evidence behind each contribution are developed in
Chapter~5. Full bibliographic details are given in \emph{Publications on PhD Work}.

\section{Research Contribution 1 --- Integrated multi-source model (RO1)}

A comprehensive multi-source model was designed, implemented and evaluated for the Indian
equity market. The Multi-Source Fusion Network of \pubcite{P1} integrates four streams ---
technical indicators from National Stock Exchange price and volume records, news sentiment
from a domain-adapted transformer encoder, social media sentiment from a hybrid lexicon and
recurrent classifier, and Indian macroeconomic series --- through modality-specific
encoders, a cross-modal attention layer that assigns time-varying relevance, and a
bidirectional recurrent encoder. It was validated on ten years of data across ten sectoral
indices under a strictly chronological split with walk-forward confirmation. The
contribution is not that fusion improves accuracy, which was already known, but that the
improvement is attributable: ablation isolates the marginal value of each stream, and the
attention distribution shows source importance shifting with the market regime.

That observation exposed a limitation in the fusion mechanism itself, addressed in
\pubcite{P4}. A Bayesian dynamic uncertainty-weighted framework models each source with its
own expert network, estimates each expert's predictive uncertainty continuously from the
variance of its recent errors, and obtains the fusion weights by a softmax over the
negative uncertainties, so that a source is discounted at the moment its predictions become
unreliable rather than at the next retraining cycle. This is the methodological core of the
objective: it converts fusion from a question of \emph{what} to combine into a question of
\emph{how} to combine under changing reliability.

\section{Research Contribution 2 --- Hybrid LLM and indicator framework (RO2)}

A hybrid framework combining large language models with conventional financial indicators
was designed and evaluated in \pubcite{P3}. News headlines are read by a transformer
sentiment encoder that produces a continuous polarity score rather than a categorical
label, preserving intensity that binary classification discards; the score is time-aligned
to trading sessions with explicit handling of closing hours and publication timestamps,
then fused at feature level with fifteen technical indicators and passed to a
gradient-boosted ensemble, against technical-only baselines under temporal
cross-validation.

Three findings constitute the contribution. The sentiment score is the highest-gain
predictor in the fused model, ahead of established technical features, so language-model
sentiment carries information price history does not. Its low correlation with price-based
features confirms that the information is orthogonal rather than redundant. And ---
reported without softening --- the incremental gain in directional accuracy is small, with
all models close to the random classifier at a single-name daily horizon; the framework's
practical value lies instead in producing the highest prediction confidence of the models
examined and the strongest performance on buy signals. Delivering it in a reproducible,
freely accessible environment was an explicit design goal, so the result can be
independently re-derived.

\section{Research Contribution 3 --- Explainable and causal framework (RO3)}

Transparency was addressed mechanistically rather than by post-hoc attribution alone. The
Dynamic-Causal Hybrid Framework of \pubcite{P2} replaces correlation-based feature
selection with a directed causal adjacency matrix recovered by conditional independence
testing, so that the features entering the model are those with demonstrable directed
influence rather than those that merely co-move; temporal structure is then captured by a
multi-scale memory encoder at intraday, daily and weekly granularity, and the signal is
converted into holdings by a credibilistic optimiser under a range directional measure
formulation. Applied to NIFTY50 constituent data over 308 trading sessions, it recovers an
interpretable causal graph, identifies a systematic causal hub, and demonstrates that a
strong directed causal influence can exist between securities whose linear correlation is
weak --- the justification for the entire approach, since correlational selection would
have discarded a genuine driver.

The second half of the objective, sentiment noise and sentiment-to-price temporal dynamics,
is addressed by \pubcite{P3} and \pubcite{P5}. The former measures accuracy separately
under positive, neutral and negative sentiment regimes and shows that sentiment spikes
frequently precede price movements while the mapping from magnitude to direction is not
simple. The latter treats the sentiment construct as a testable claim rather than an
assumption: the composite is validated against four independently published series, carries
the expected sign against every one in every sub-period, and separates the identified
market states at $1.8 \times 10^{-8}$ --- and the same study then measures separately
whether the validated composite forecasts anything, and reports that it does not. Keeping
construct validity and predictive usefulness apart is itself a contribution, because
conflating them is a common error in this literature.

\section{Research Contribution 4 --- Validated decision framework and cross-market
transfer (RO4)}

A decision-theoretic framework converting a fused signal into an executable portfolio was
specified and evaluated in \pubcite{P5}. Latent market states are identified by a Gaussian
mixture refitted at every rebalance with explicit label anchoring; composites are fused
with state-dependent weights and scaled onto a return metric by information-ratio scaling;
the scaled signal enters a constrained quadratic program with covariance shrinkage,
position caps and a turnover penalty; and the sleeve is scaled by a regime-dependent risk
budget. The design deliberately separates what to hold relative to everything else from how
much total risk to carry, so that each layer can be removed and tested independently.
Evaluation spans eighteen years of daily data on ten multi-asset instruments with
transaction costs, a fully reported specification sweep, four stress episodes selected on
documented market history, a 665-session holdout postdating every design decision, and a
secondary mega-capitalisation universe as a cross-market test.

The contribution has three parts. The first is the framework itself, specified in enough
detail to be reproduced from the description alone. The second is an ablation design that
isolates component contributions at matched average exposure --- necessary because scaling
a portfolio toward cash changes return and volatility together and would otherwise confound
risk timing with risk level. The third, most useful to a practitioner, is the
identification of two failure modes: an absolute volatility target never reached on a
diversified universe, which silently disables the regime layer, and a signal-to-allocation
conversion in which mismatched units rather than information determine the optimiser's
solution. Both were found only by disassembling the framework, and both generalise beyond
it. The cross-market component is answered honestly: transferred unchanged to a
concentrated universe the framework retains its defensive behaviour but loses on
risk-adjusted return, for reasons that are a property of that universe rather than of the
framework --- a bounded claim, which is more useful than an unbounded one.

\section{Consolidated Mapping of Published Work to Objectives}

Table~\ref{tab:mapping} states the relationship between each publication and each
objective. Two conventions are used throughout the report. A \emph{primary} mapping means
that the publication was designed to answer the objective and supplies its principal
evidence. A \emph{partial} mapping means that the publication contributes a specific,
identifiable component of the objective without answering it as a whole. Where no
defensible relationship exists, none is claimed; the blank cells in
Table~\ref{tab:mapping} are deliberate and should be read as an accurate statement of
scope rather than as an omission.

\begin{table}[!htb]
\centering
\caption[Mapping of published work to the research objectives]{Mapping of published work to the research objectives, with the nature of the
contribution in each case}
\label{tab:mapping}
\tabfont
\begin{tabular}{@{}L{1.1cm}L{3.5cm}cccc L{4.4cm}@{}}
\toprule
\textbf{Ref.} & \textbf{Publication (short title)} & \textbf{RO1} & \textbf{RO2} &
\textbf{RO3} & \textbf{RO4} & \textbf{Nature of the contribution} \\
\midrule
\pubcite{P1} & Multi-source data fusion framework (conference) & \textbf{P} & p & --- & p
& Four-stream architecture, cross-modal attention, source ablation, ten-sector
evaluation \\
\pubcite{P2} & Dynamic-causal hybrid framework, NIFTY50 (conference) & p & --- &
\textbf{P} & p & Directed causal feature selection, multi-granularity memory,
credibilistic allocation \\
\pubcite{P3} & LLM-extracted news sentiment and market data (journal) & p & \textbf{P} & p
& --- & LLM sentiment pipeline, feature-level fusion, sentiment-regime accuracy analysis \\
\pubcite{P4} & Bayesian dynamic uncertainty-weighted fusion (journal) & \textbf{P} & p &
p & --- & Time-varying source reliability, softmax weighting over negative uncertainties \\
\pubcite{P5} & Signal fusion to asset allocation (journal) & --- & --- & p & \textbf{P} &
Regime detection, constrained allocation, matched-exposure ablation, holdout and
cross-market test \\
\pubcite{P6} & Deep learning in stock forecasting (conference) & p & p & --- & --- &
Comparative groundwork that established the research gap \\
\bottomrule
\multicolumn{7}{@{}l@{}}{\scriptsize \textbf{P} = primary contribution;\quad p = partial
contribution;\quad --- = no contribution claimed.}
\end{tabular}
\end{table}

Figure~\ref{fig:contribmap} shows the same relationship as a flow, which makes the
structure of the research programme visible: two publications feed the integration layer,
one supplies the language-model signal, one supplies transparency and one supplies the
decision layer, and the evidence accumulates from left to right rather than in parallel.

\begin{figure}[!htb]
\centering
\scalebox{0.86}{\begin{tikzpicture}[node distance=4mm]
\node[blk,minimum width=31mm,minimum height=7.5mm] (p1) {\pubcite{P1} MSFN\\four-stream fusion};
\node[blk,below=2.6mm of p1,minimum width=31mm,minimum height=7.5mm] (p4) {\pubcite{P4} Bayesian\\uncertainty fusion};
\node[blk,below=2.6mm of p4,minimum width=31mm,minimum height=7.5mm] (p3) {\pubcite{P3} LLM sentiment\\ + indicators};
\node[blk,below=2.6mm of p3,minimum width=31mm,minimum height=7.5mm] (p2) {\pubcite{P2} DCHF\\causal + multi-scale};
\node[blk,below=2.6mm of p2,minimum width=31mm,minimum height=7.5mm] (p5) {\pubcite{P5} decision-theoretic\\allocation};
\node[blkd,below=2.6mm of p5,minimum width=31mm,minimum height=7mm] (p6) {\pubcite{P6} comparative\\groundwork};

\node[blkg,right=30mm of p1,minimum width=26mm,minimum height=7mm] (r1) {\textbf{RO1}\\integration};
\node[blkg,below=6.5mm of r1,minimum width=26mm,minimum height=7mm] (r2) {\textbf{RO2}\\hybrid LLM};
\node[blkg,below=6.5mm of r2,minimum width=26mm,minimum height=7mm] (r3) {\textbf{RO3}\\transparency};
\node[blkg,below=6.5mm of r3,minimum width=26mm,minimum height=7mm] (r4) {\textbf{RO4}\\validation};

\node[blka,right=20mm of r2,minimum width=26mm,minimum height=22mm,yshift=-8mm]
  (th) {\textbf{Thesis}\\[2pt] integrated, transparent\\ and validated\\ decision framework};

\draw[ar] (p1) -- (r1);
\draw[ar] (p4) -- (r1);
\draw[arb] (p1.east) to[out=0,in=180] (r2.west);
\draw[ar] (p3) -- (r2);
\draw[arb] (p3.east) to[out=0,in=180] (r3.west);
\draw[ar] (p2) -- (r3);
\draw[arb] (p2.east) to[out=0,in=180] (r4.west);
\draw[ar] (p5) -- (r4);
\draw[arb] (p5.east) to[out=0,in=185] (r3.west);
\draw[arb] (p6.east) to[out=0,in=190] (r1.west);
\draw[ar] (r1.east) to[out=0,in=180] (th.west);
\draw[ar] (r2.east) to[out=0,in=180] (th.west);
\draw[ar] (r3.east) to[out=0,in=180] (th.west);
\draw[ar] (r4.east) to[out=0,in=180] (th.west);
\node[lbl,above=1.5mm of p1] {Published work};
\node[lbl,above=1.5mm of r1] {Objectives};
\end{tikzpicture}}
\caption[Contribution map of published work to objectives]{Contribution map. Solid arrows denote primary contributions and dashed arrows
denote partial contributions, matching Table~\ref{tab:mapping}.}
\label{fig:contribmap}
\end{figure}
